% Auto-generated by tools/qbe.py project-article-update.
% Do not edit this file by hand; edit the proof artifacts or article sections instead.

\section{Generated Current Cycle Status}
\label{app:generated-cycle-status}

This appendix page is generated from the latest ABEIS proof cycle.  It is a
status bridge for article maintenance, not a polished theorem proof.  Stable
mathematical claims should be copied into the main text only after the
corresponding Lean declaration, source-paper anchor, and proof-note export are
available.

\paragraph{Cycle.}
Task \texttt{QBE-AUTO-002}, cycle \texttt{1}, run
\texttt{runs/20260611-214323-QBE-AUTO-002-cycle01}.  Generated at
\texttt{2026-06-11 22:00:56}.

\paragraph{Lean status signal.}
\begin{itemize}
  \item \texttt{QuantumBlockEncoding/RobinMatrix.lean:24024:  sorry}
  \item \texttt{QuantumBlockEncoding/RobinMatrix.lean:24062:  sorry}
\end{itemize}

\paragraph{Plain-language status.}
The current GHL case study is not blocked because the paper lacks a proof sketch.  It is blocked because ABEIS must turn the circuit in the paper into an exact matrix statement and prove that the clean block entry has the coefficient claimed by the paper.  In ordinary terms, the paper says that a specific path through the circuit leaves the desired coefficient, while Lean requires the corresponding matrix entry and all unwanted branches to be proved exactly.

\paragraph{Pre-Lean verifier candidates.}
These checks are necessary-condition filters, not proofs.  A failed exact
finite check usually indicates that the current Lean target, circuit
transcript, or index map is wrong.  A passing check only means the candidate
survived a cheaper diagnostic; the final claim still requires a named Lean
theorem or an explicit contract.
\begin{itemize}
  \item \textbf{active [0,0] entry.} Fast check: exact rational matrix or path-sum evaluation.  Necessary because if the exact finite entry is not the target coefficient, the Lean equality for that entry cannot be true.  Lean still must prove: a passing check is only a counterexample filter; Lean must still prove the named entry lemma.
  \item \textbf{remaining branch vanish/cancel.} Fast check: support and path checker for the remaining backend slots.  Necessary because if an unwanted clean-branch path survives numerically, the block projection cannot match the paper target.  Lean still must prove: Lean must still prove the zero/cancellation lemma in the formal circuit semantics.
  \item \textbf{Ry boundary convention.} Fast check: symbolic 2-by-2 rotation check.  Necessary because a mismatched half-angle convention changes the boundary amplitude before any Lean tactic is relevant.  Lean still must prove: Lean must still record the convention bridge as a source-supported theorem.
  \item \textbf{sparse-access map.} Fast check: finite range/injectivity/permutation check.  Necessary because a reversible oracle cannot exist for a colliding or out-of-range finite map.  Lean still must prove: Lean must still prove the reversible extension and cleanup obligations.
  \item \textbf{coefficient oracle clean branch.} Fast check: exact finite evaluator for f(x\_i)/N\_f.  Necessary because the final block entry uses this coefficient, so a wrong clean branch invalidates the target theorem.  Lean still must prove: Lean must still prove bounds, orthogonality, and unitary completion or keep them as contracts.
\end{itemize}

\paragraph{Current proof-DAG frontier.}
\begin{itemize}
  \item active\_selected\_slot\_eval\_comparison\_leaf: evaluated seven-gate \textasciigrave{}[0,0]\textasciigrave{} entry equals the selected slot-\textasciigrave{}2\textasciigrave{} contribution; status: active mathematical leaf; Lean: target \textasciigrave{}ActiveSelectedSlotEvalComparison\textasciigrave{} displayed above
  \item active\_uncast\_to\_prepared\_entry\_leaf: uncast seven-gate \textasciigrave{}[0,0]\textasciigrave{} entry equals the cached prepared entry; status: open dependent target; Lean: target \textasciigrave{}ActiveUncastToPreparedEntry\textasciigrave{}
  \item source\_prepared\_entry\_leaf: \textasciigrave{}SourcePreparedEntry\textasciigrave{}: active/prepared entry equality; status: open dependent target; Lean: \textasciigrave{}(oneTermRobinGamma3BoundaryActivePreparedEntryTarget\_n3 H).entryEqualityStatement\textasciigrave{}
  \item unitary\_fold\_leaf: \textasciigrave{}FullUnitaryFold\textasciigrave{}: full signal-zero unitary entry equals the seven-slot backend branch fold; status: open dependent root; Lean: target \textasciigrave{}oneTermRobinGamma3BoundaryProjectionSummationTarget\_n3.signalUnitaryEntry = blockExtractionBranchContributionSum oneTermRobinGamma3BoundaryBackendBranchContribution\_n3\textasciigrave{}
  \item backend\_expansion\_leaf: backend-expansion statement for the branch-contribution target; status: open equivalent endpoint; Lean: \textasciigrave{}oneTermRobinGamma3BoundaryBackendBranchContributionTarget\_n3.backendExpansionStatement\textasciigrave{}; \textasciigrave{}oneTermRobinGamma3BoundaryBackendExpansionStatement\_equivUnitaryEntryFold\_n3\textasciigrave{}
\end{itemize}

\paragraph{Open obligation signal.}
\begin{itemize}
  \item theorem-facing Fig. 4 transcript: Lean \textasciigrave{}GHL2025.oneTermRobinTheoremFacingFig4Circuit\_gateList\textasciigrave{}; \textasciigrave{}GHL2025.oneTermRobinGate\_U\_indic\_dagger\_selfInverseBridge\textasciigrave{}; class GHL-internal transcript plus QBE-local indicator bridge; status compiled guard; no full semantic proof
  \item prepared entry backend-fold normal form: Lean \textasciigrave{}oneTermRobinGamma3BoundaryActivePreparedEntryTarget\_preparedEntry\_eq\_backendFold\_n3 H hUniform\textasciigrave{}; class QBE-local prepared-side semantic bridge under external \textasciigrave{}HUniform\textasciigrave{} contract; status proved feeder; retired as next lower target
  \item active wrapper/cast removal: Lean \textasciigrave{}oneTermRobinGamma3BoundaryActivePreparedEntryTarget\_iff\_uncastActiveEntry\_n3 H\textasciigrave{}; class QBE-local finite matrix-entry bridge; status proved feeder; retired as next lower target
  \item nonselected backend-slot evaluated vanish family: Lean slots \textasciigrave{}0\textasciigrave{}, \textasciigrave{}1\textasciigrave{}, \textasciigrave{}3\textasciigrave{}, \textasciigrave{}4\textasciigrave{}, \textasciigrave{}5\textasciigrave{}, and \textasciigrave{}6\textasciigrave{} evaluated branch vanish theorems; class QBE-local finite branch evaluation lemmas; status proved feeders; retired as next lower targets
  \item selected slot-\textasciigrave{}2\textasciigrave{} backend-fold reduction: Lean \textasciigrave{}oneTermRobinGamma3BoundaryBackendBranchFoldEval\_eq\_selectedSlotContribution\_n3 env\textasciigrave{}; class QBE-local finite branch evaluation lemma; status proved feeder; retired as next lower target
  \item active-selected comparison to evaluated-fold bridge: Lean \textasciigrave{}oneTermRobinGamma3BoundaryActiveSelectedSlotEvalComparison\_iff\_evaluatedBackendFold\_n3 env\textasciigrave{}; class QBE-local route-packaging bridge; status proved bridge; retired as next lower target
  \item active selected-slot evaluated comparison: Lean target \textasciigrave{}ActiveSelectedSlotEvalComparison\textasciigrave{}; class QBE-local finite seven-gate matrix semantics; status active mathematical leaf
  \item prepared-sandwich semantics fallback: Lean \textasciigrave{}oneTermRobinGamma3BoundaryPreparedCircuitSemanticsGap\_n3 H\textasciigrave{}; \textasciigrave{}oneTermRobinGamma3BoundaryUncastPreparedSandwichEvalStatement\_n3 H env\textasciigrave{}; class QBE-local source-prepared matrix semantics; status controlled fallback if the direct H-free comparison produces a shape/register gap
\end{itemize}

\paragraph{Open GHL contribution obligations.}
\begin{itemize}
  \item \texttt{Uindic}: Indicator unitary; status Permutation and self-inverse helpers compile, but this is not yet the final block-encoding proof..
  \item \texttt{RyBoundary}: Boundary-controlled Ry rotations; status Active convention audit: the Ry angle convention must be source-supported before closing the boundary amplitude proof..
  \item \texttt{RobinTheorem}: One-term Robin block-encoding theorem; status Main active target: the theorem-facing Lean statement is not closed..
  \item \texttt{GammaSlices}: Gamma wavefunction slices and coefficient product; status Several finite boundary lemmas exist; the final projection/product bridge is still open..
  \item \texttt{FigRobin}: Figure 4 Robin circuit transcript; status Transcript guard compiles for visible indicator dagger and H preparation; the active seven-gate backend remains a component..
  \item \texttt{OneD}: One-dimensional Hamiltonian block encoding; status Deferred until the one-term Robin theorem closes..
  \item \texttt{MultiD}: Multidimensional block encoding; status Planned; depends on one-term and LCU abstractions..
\end{itemize}

\paragraph{Open external technical lemma obligations.}
\begin{itemize}
  \item \texttt{tl-ghl-lemma1-banded-sparse-access}: contract-only; next action Keep as explicit external contract until the exact cited construction is formalized or imported as a theorem..
  \item \texttt{tl-ghl-lemma3-sparse-amplitude}: contract-only; next action Maintain the clean-branch contract and isolate any missing unitarity proof as an external technical lemma..
  \item \texttt{tl-ghl-theorem5-piecewise-polynomial-of}: contract-only; next action Use only as a named contract in the current theorem; do not invent stronger smoothness or range assumptions..
  \item \texttt{tl-uniform-sparse-register-preparation}: contract-only; next action Keep as a theorem-facing contract unless the exact state-preparation proof is needed to close a resource theorem..
  \item \texttt{tl-ry-boundary-amplitude-convention}: obligation; next action Do not silently change the paper angle; prove the convention bridge or record the exact cited convention..
  \item \texttt{tl-qsvt-blockencoding-simulation}: paper-cited; next action Defer until the one-term Robin theorem and LCU combination are closed..
  \item \texttt{tl-clean-block-definition}: contract-only; next action Use as the final target predicate; ensure circuit entry lemmas are bridged to this semantic definition..
\end{itemize}

\paragraph{Recent typed verifier feedback.}
\begin{itemize}
  \item Leaf \texttt{active\_selected\_slot\_eval\_comparison}, class \texttt{shape\_or\_register\_gap}; next route: reroute to prepared-sandwich semantics gap unless a direct source-faithful evalGateMatrices path proves ActiveSelectedSlotEvalComparison.
  \item Leaf \texttt{active\_selected\_slot\_eval\_comparison\_bridge}, class \texttt{None}; next route: prove ActiveSelectedSlotEvalComparison itself, then recover evaluated backend fold through oneTermRobinGamma3BoundaryActiveSelectedSlotEvalComparison\_iff\_evaluatedBackendFold\_n3.
  \item Leaf \texttt{post\_slot6\_active\_uncast\_source\_prepared\_design\_concurrency\_note}, class \texttt{symbolic\_bridge\_gap}; next route: compare active uncast evalGateMatrices [0,0] entry with selected slot2 contribution.
  \item Leaf \texttt{post\_slot6\_active\_uncast\_source\_prepared\_design}, class \texttt{symbolic\_bridge\_gap}; next route: prove source-prepared evaluated active leaf or strict selected-slot/active-side evaluated feeder.
  \item Leaf \texttt{selected\_slot2\_backend\_fold\_eval\_feeder}, class \texttt{None}; next route: compare active uncast evalGateMatrices [0,0] entry to selected slot-2 contribution to close ActiveUncastToPreparedEntry/evaluated active-prepared field.
\end{itemize}

\paragraph{Recent proof-attempt memory.}
\begin{itemize}
  \item \texttt{proof-attempts/QBE-AUTO-002/unitary-fold-leaf-dag-addendum-20260610-1156-lower1.md}
  \item \texttt{proof-attempts/QBE-AUTO-002/post-active-selected-slot-bridge-middle-packet-20260611-2148.md}
  \item \texttt{proof-attempts/QBE-AUTO-002/post-selected-slot-active-uncast-middle-packet-20260611-2129.md}
  \item \texttt{proof-attempts/QBE-AUTO-002/post-slot-six-active-uncast-middle-packet-20260611-2038.md}
  \item \texttt{proof-attempts/QBE-AUTO-002/post-slot-five-evaluated-remaining-slots-middle-packet-20260611-2018.md}
\end{itemize}

\paragraph{Article-writing rule.}
The project report may discuss this cycle as process evidence.  It must not
claim completion of the first case study \citep{guseynovHuangLiu2025} unless
the final theorem-facing block-extraction statement is closed in Lean and the
Markdown/LaTeX proof map has been synchronized.
